\begin{frame}[fragile]{실습: 완전탐색이 가능한 경우(틱택토)}
  ``모델을 안다''고 해서 항상 완전탐색(minimax)을 쓸 수 있는 건
  아니다 --- \textbf{상태 수가 관건}. 틀만 같고 상태 수가 훨씬
  작은 틱택토(3$\times$3)에서 검증:
\begin{lstlisting}
WIN_LINES = [(0,1,2),(3,4,5),(6,7,8),(0,3,6),(1,4,7),(2,5,8),(0,4,8),(2,4,6)]
nodes_visited = 0

def winner(board):
    for a, b, c in WIN_LINES:
        if board[a] != 0 and board[a] == board[b] == board[c]:
            return board[a]
    return 0

def minimax(board, player):
    global nodes_visited
    nodes_visited += 1
    w = winner(board)
    if w != 0: return w
    if all(x != 0 for x in board): return 0   # 무승부
    scores = []
    for i in range(9):
        if board[i] == 0:
            board[i] = player
            scores.append(minimax(board, -player))
            board[i] = 0
    return max(scores) if player == 1 else min(scores)

result = minimax([0]*9, player=1)
print("완전탐색 결과:", result)
print("탐색한 전체 국면 수:", nodes_visited)
\end{lstlisting}
\end{frame}
